% Generated draft: requires mathematical and source review.
\documentclass[11pt]{article}
\usepackage[margin=25mm]{geometry}
\usepackage{fontspec}
\setmainfont{Latin Modern Roman}
\usepackage{amsmath,amssymb,mathtools,unicode-math}
\setmathfont{Latin Modern Math}
\usepackage{enumitem,needspace,xurl,hyperref}
\setlength{\parindent}{0pt}
\setlength{\parskip}{5pt}
\setlength{\emergencystretch}{3em}
\newcommand{\N}{\mathbb{N}}
\newcommand{\Z}{\mathbb{Z}}
\newcommand{\Q}{\mathbb{Q}}
\newcommand{\R}{\mathbb{R}}
\newcommand{\C}{\mathbb{C}}
\newcommand{\F}{\mathbb{F}}
\newcommand{\A}{\mathbb{A}}
\newcommand{\PP}{\mathbb{P}}
\newcommand{\E}{\mathbb{E}}
\newcommand{\eps}{\varepsilon}
\newcommand{\dd}{\,\mathrm{d}}
\begin{document}
\section*{UnsolvedMath ID: 30001047}
\textbf{Source problem code:} catalogue entry
\subsection{Definitions and mathematical statement}
Let \(M\) be a hyperbolic tunnel-number-one \(3\)-manifold, and let \(\tau\) be an unknotting tunnel in the sense of the cited report. In the familiar one-cusped setting, one works with a compact core \(\overline M\) and a properly embedded arc \(\tau\) with endpoints on its boundary, such that deleting a regular neighborhood of the arc leaves a handlebody. The usual knot-exterior case has torus boundary and genus-two handlebody complement. The record does not state that all its manifolds are restricted to this setting, so this is an explanatory model rather than an added hypothesis.

The first question asks whether every such tunnel is isotopic to a geodesic representative in the given hyperbolic metric, with the endpoint or cusp convention intended in the report. Isotopy is part of the assertion; it is not replaced by the weaker existence of a geodesic in a homotopy class, and the representative must have the embeddedness needed for an isotopic arc.

The second question asks whether unknotting tunnels have bounded length. The substantive uniform reading is whether a constant \(L<\infty\), independent of the allowed manifold and tunnel, bounds the intended tunnel-length invariant. The source does not specify whether this invariant is the infimum over representatives, the length of a geodesic representative, or a truncated length relative to chosen cusp neighborhoods. In a cusped setting, an arc with ideal endpoints has infinite full hyperbolic length, so a finite-length question needs such a convention. Neither a truncation nor a weaker manifold-dependent bound is silently selected.

The supplied assessment reports the uniform bounded-length assertion false, with tunnel-number-one examples having arbitrarily long tunnels, while reporting the general geodesic-isotopy question open. These component statuses are preserved as supplied claims rather than combining the two questions into one open conjecture.
\subsection{Short English statement}
Must every unknotting tunnel in a hyperbolic tunnel-number-one manifold be isotopic to a geodesic, and is there a uniform bound for the intended tunnel length? The source reports the length bound false and the geodesic question open, subject to source verification.
\subsection{Sources}
\begin{itemize}
\item Local UnsolvedMath registry, permanent ID \(30001047\), source code catalogue entry, ``Geodesic Representatives and Lengths of Unknotting Tunnels.'' No individual proposer is supplied.
\item \emph{Topologie} (2009), Oberwolfach Reports, DOI 10.4171/owr/2008/43, pp. 2419--2476: \url{https://doi.org/10.4171/owr/2008/43}.
\item Supplied reference, ``The length of unknotting tunnels'': \url{https://arxiv.org/abs/0812.0858}.
\item Supplied reference, ``Geodesic representatives of tunnels'': \url{https://arxiv.org/abs/1302.5469}.
\item Supplied literature triage, checked 2026-08-21: status ``partially solved.'' It reports unbounded tunnel lengths in hyperbolic tunnel-number-one examples, disproving the bounded-length assertion, and says whether every tunnel has a geodesic representative remains open. This is supplied dated curation, not independently verified; time-sensitive status, length conventions, and cited claims require independent verification.
\end{itemize}
\end{document}
